\section{Method}
\label{sec:method}

We study \emph{ParT-EV}, a family of Particle Transformer extensions with
edge-conditioned value messages. We use \emph{EdgeValue} for the message
mechanism and \emph{ParT-EV} for models that include it. The relation feature
set and the choice of shared or layerwise bias are specified separately;
the baseline and bias-only ablations are comparison models, not ParT-EV variants.

\subsection{Particle attention and pair information}

The Particle Transformer (ParT) backbone~\cite{Qu2022ParticleTransformer} has a
particle embedding with widths $(128,512,128)$, eight particle self-attention
blocks with eight heads, and two class-attention blocks followed by the
classifier. The attention width is $128$, giving head dimension $d_h=16$.
Our question is whether additional relation descriptors, layer-dependent
interaction biases, or relation-conditioned value messages improve this baseline.

For a valid query particle $i$ and context particle $j$, the baseline supplies four kinematic pair features,
\begin{equation}
 r^{\mathrm{std}}_{ij}
 = \bigl(\Lambda(k_{t,ij}),\Lambda(\zeta_{ij}),
          \Lambda(\Delta_{ij}),\Lambda(m^2_{ij})\bigr),
 \label{eq:standard-pairs}
\end{equation}
where $\Lambda(x)=\log(\max\{x,10^{-8}\})$. Here $\Delta$ is separation in
rapidity--azimuth space, $k_t$ the relative transverse-momentum scale,
$\zeta$ the smaller transverse-momentum share, and $m^2$ the pair invariant
mass squared. The rapidity used here is distinct from the pseudorapidity
used in the TRACK and REGION descriptors below. The baseline therefore already
incorporates pairwise kinematics within global attention.

Writing $\ell$ for the layer and $h$ for the head, biased scaled dot-product attention~\cite{Vaswani2017Attention,Qu2022ParticleTransformer} is
\begin{equation}
 \begin{aligned}
 L^{\ell h}_{ij}
 &=\frac{(q_i^{\ell h})^{\mathsf T}k_j^{\ell h}}{\sqrt{d_h}}
       +b^{\ell h}_{ij},\\
 a^{\ell h}_{ij}
 &=\frac{\exp L^{\ell h}_{ij}}
        {\sum_{k\in\mathcal V}\exp L^{\ell h}_{ik}},
 \qquad i,j\in\mathcal V.
 \end{aligned}
 \label{eq:biased-attention}
\end{equation}
Here $\mathcal V$ is the set of valid constituents, so the normalization excludes
padded keys and permits self-attention ($j=i$). We write
$\widetilde a=\operatorname{Dropout}(a)$ for the attention weights after
training-time dropout; at evaluation $\widetilde a=a$.
Bias-only attention aggregates $\sum_{j\in\mathcal V}\widetilde a^{\ell h}_{ij}v_j^{\ell h}$:
relations affect the weights, while each context value remains independent of its query.

\subsection{Selected relation descriptors}

The selected input concatenates the standard four channels with learned PT, TRACK, and REGION embeddings of dimensions $8$, $12$, and $12$, respectively, giving $r_{ij}\in\mathbb R^{36}$. Throughout, $i$ denotes the query and $j$ the context, and directional features are not symmetrized.

\paragraph{PT.}
Let $S_T=\sum_{k\in\mathcal V}p_{T,k}$ be the scalar constituent-$p_T$ sum
and $\epsilon=10^{-6}$ the numerical safeguard used by the added relation
families. Each endpoint supplies $p_{T,i}/(S_T+\epsilon)$, the safeguarded
log-fraction $\log[(p_{T,i}+\epsilon)/(S_T+\epsilon)]$, and its descending
normalized $p_T$ rank. The remaining four channels are the context-minus-query
log-fraction and rank differences, the pair log-fraction
$\log[(p_{T,i}+p_{T,j})/(S_T+\epsilon)+\epsilon]$, and
$(p_{T,j}-p_{T,i})/(p_{T,j}+p_{T,i}+\epsilon)$. Tied ranks are averaged.
A $10\!\to\!32\!\to\!8$ network embeds these ten quantities.

\paragraph{TRACK.}
Each endpoint supplies $d_0$, $d_z$, logarithmic effective uncertainties, inverse-hyperbolic-sine impact-parameter significances, and a validity flag. Validity requires charged-hadron, electron, or muon identification, finite, nonmissing measurements, and positive uncertainties. Effective errors satisfy $\sigma_{u,\mathrm{eff}}^2=\sigma_u^2+f_u^2$ for $u\in\{d_0,d_z\}$, with floors $f_u$ fitted on training data. Shared endpoint weights produce $g_i,g_j\in\mathbb R^{16}$ using a $7\!\to\!32\!\to\!16$ network. The pair encoder receives
\begin{equation}
 [g_i,g_j,g_j-g_i,g_i\odot g_j,s_{ij},c_{ij}]\in\mathbb R^{81},
 \label{eq:track-pairs}
\end{equation}
where $s_{ij}$ is a four-component one-hot encoding of the ordered endpoint
validity states and $c_{ij}$ contains 13 compatibility channels. These comprise
uncertainty-normalized context-minus-query displacement differences, their
squared-sum $\chi^2$ through $\log(1+\chi^2)$ and
$\exp[-\min(\chi^2,25)/2]$, extrema of absolute transformed significances and
their signed products, and angular descriptors $\sin\Delta\phi$,
$\cos\Delta\phi$, and $\log(\Delta R+\epsilon)$, with
$\Delta\phi=\phi_i-\phi_j$ wrapped periodically and $\Delta R$ measured in
pseudorapidity--azimuth space. Compatibility channels vanish unless both
tracks are valid. The resulting $81\!\to\!48\!\to\!12$ network retains
missing-measurement information through the validity states.

\paragraph{REGION.}
A deterministic, beam-free angular tree repeatedly merges the closest pair in $(\eta,\phi)$ using four-vector addition. It follows the angular-clustering motivation of Cambridge/Aachen~\cite{Dokshitzer1997Cambridge,Wobisch1999Cambridge}, with the explicit pseudorapidity-based prescription above. Exclusive partitions at $K=2,4,8$ clusters, capped by constituent multiplicity, provide 41 features: three same-cluster indicators; five lowest-common-ancestor descriptors (normalized depth, merge angle, $k_t$, momentum sharing, and mass); 18 endpoint cluster $p_T$, mass, and multiplicity fractions; six constituent shares within those clusters; six constituent-to-cluster-axis distances; and three context-minus-query cluster-$p_T$ rank differences. Angles use reference radius $R_{\mathrm{ref}}=0.8$, merge $k_t$ uses $p_{T,\mathrm{jet}}R_{\mathrm{ref}}$, and mass uses jet mass. Cluster $p_T$ and mass fractions, and merge angle, $k_t$, and mass descriptors use logarithms. Undefined diagonal merge quantities are zeroed. A $41\!\to\!32\!\to\!12$ network embeds the result. These descriptors are supplied for every valid pair, while attention remains global.

Here jet $p_T$ and mass are obtained from the summed constituent four-vector,
not the scalar sum $S_T$ used by PT. Ancestor depth is divided by the maximum
leaf depth, with the denominator bounded below by one.

All family networks use GELU and RMS normalization between their two linear
maps and are trained jointly with the classifier. Before entering these
networks, designated continuous features are centered and rescaled using
statistics from a fixed subset of 50,000 training jets. This subset is used
only to estimate preprocessing constants; the classifier is trained on the
full one-million-jet training sample. For each feature, we subtract its median,
divide by a robust measure of its spread, and clip extreme values:
\begin{equation}
 \widehat x=\operatorname{clip}_{[-8,8]}
 \left(\frac{x-\operatorname{median}(x)}
 {\max[(Q_{75}-Q_{25})/1.349,10^{-6}]}\right).
 \label{eq:relation-normalization}
\end{equation}
Here $Q_{25}$ and $Q_{75}$ are the training-sample quartiles; the resulting
scale is less sensitive to outliers than a standard deviation. Track
uncertainty floors are also estimated from training data. Binary and
designated fixed-scale features are left unchanged. These preprocessing
constants remain fixed during validation and testing, and padded pair
embeddings are zeroed.

\subsection{Shared and layerwise interaction biases}

A three-stage, width-64 pair network maps either input set to a common relation stem $z_{ij}=\phi(r_{ij})\in\mathbb R^{64}$. Shared bias uses one head-output map $\beta$ in all particle layers, $b^{\ell h}_{ij}=\beta_h(z_{ij})$. Layerwise bias uses independently trainable maps $\beta_\ell$, so $b^{\ell h}_{ij}=\beta_{\ell h}(z_{ij})$, while retaining the common stem. The layerwise maps, including their output normalization, have identical initial parameters.

\subsection{Relation-conditioned value messages}
\label{sec:edgevalue}

Pairwise attention bias and pairwise message content play different roles.
In bias-only ParT, the relation between particles $i$ and $j$ affects the scalar
attention weight, but the value vector $v_j^{\ell h}$ is the same for every
query receiving information from $j$ within that layer and head. EdgeValue
provides an additional route for relation information to enter the particle
representation: it augments the value with a vector that depends on both
endpoints. The same context particle can therefore contribute different message
content to different queries, rather than differing only in its attention weight.
Building on relation-aware value messages~\cite{Shaw2018Relative}, ParT-EV uses
\begin{align}
 y_i^{\ell h}
 &=\sum_{j\in\mathcal V}\widetilde a^{\ell h}_{ij}
   \left(v_j^{\ell h}+U_{\ell h}z_{ij}\right),
 \label{eq:edgevalue}\\
 \delta_i^{\ell h}
 &=\sum_{j\in\mathcal V}\widetilde a^{\ell h}_{ij}U_{\ell h}z_{ij},
 \qquad U_{\ell h}\in\mathbb R^{16\times64}.
 \label{eq:edgevalue-message}
\end{align}
Here $y_i^{\ell h}$ is the head aggregate before the existing output projection,
and $\delta_i^{\ell h}$ is the addition to the ordinary weighted value sum.
Both terms use the same attention weights $\widetilde a$ and the same
training-time attention-dropout realization.

The shared relation stem $z_{ij}=\phi(r_{ij})$ supplies both the bias and
value-message branches. Encoded from the input pair descriptors, it is shared
across the particle-attention layers rather than updated from their evolving
particle states. Its encoder is trained jointly with the classifier.
Each layer and head has its own projection $U_{\ell h}$, so this common representation can contribute
differently at different depths and through different heads. The mechanism
applies both to the standard four-feature input and to the selected relation
input, and does not require layerwise attention bias. For example, an encoded
track-compatibility or regional relationship can influence both the weight and
the vector content of a message.

The value-message projections are bias-free. Within each layer, their stacked
$8\times16\times64$ parameter tensor is initialized with Xavier-uniform
initialization and optimized jointly with the rest of the network. Setting all
$U_{\ell h}=0$ removes the additional term and recovers the corresponding
bias-only attention operation for fixed remaining parameters. Invalid pair
stems and padded-query edge contributions are zeroed. The head outputs are
concatenated and passed through the existing output projection, with the
surrounding transformer block unchanged. EdgeValue is applied to all eight
particle self-attention blocks, while the two class-attention blocks remain
unchanged. Its projections add $8\times8\times16\times64=65{,}536$ trainable
parameters relative to the matching bias-only path, separately from any
parameters added by selected relation encoders or layerwise bias projections.

Because $U_{\ell h}$ is linear and independent of $i$ and $j$, the relation
aggregation can precede projection:
\begin{equation}
 \delta_i^{\ell h}
 =U_{\ell h}\left(\sum_{j\in\mathcal V}
   \widetilde a^{\ell h}_{ij}z_{ij}\right).
 \label{eq:edgevalue-efficient}
\end{equation}
This rearrangement avoids explicit per-head pair-value vectors, while pair
features and attention weights still require quadratic storage.
